% Originally: the \section{Solutions} of content/week-02.tex and content/week-03.tex

\section{Solutions}
\label{sec:continuity_solutions}

% Transition: Claude Opus 5 (High)
Not every solution in this chapter is collected here. \Cref{ex:distance_to_point_continuous}
(\cpageref{ex:distance_to_point_continuous}) and \cref{ex:evaluation_map_continuous}
(\cpageref{ex:evaluation_map_continuous}) are solved where they appear.

\begin{exercisesolution}[Solution to \cref{ex:2.6}]
% Generator: Claude Opus 5 (High)
\textbf{$d$ is $1$-Lipschitz with respect to $d_2$.} Let $x = (x_1,x_2)$ and $y = (y_1,y_2)$ be
points of $X\times X$. Applying the triangle inequality in $X$ twice,
\[d(x_1,x_2) \ \leq\ d(x_1,y_1) + d(y_1,y_2) + d(y_2,x_2),\]
so $d(x_1,x_2) - d(y_1,y_2) \leq d(x_1,y_1) + d(x_2,y_2) = d_2(x,y)$. Exchanging the roles of
$x$ and $y$ bounds the opposite difference in the same way, hence
\[\big|d(x_1,x_2) - d(y_1,y_2)\big| \ \leq\ d_2(x,y),\]
which is $1$-Lipschitz continuity (\cref{def:lipschitz_continuous}) and therefore continuity
(\cref{rem:continuity_hierarchy}).

\textbf{For $d_1$ and $d_3$: yes, still Lipschitz, but the constant changes.} By
\cref{ex:2.5} part 2 we have $d_2 \leq 2d_1$ and $d_2 \leq \sqrt2\,d_3$, so
\[\big|d(x_1,x_2)-d(y_1,y_2)\big| \leq 2\,d_1(x,y)
  \qquad\text{and}\qquad
  \big|d(x_1,x_2)-d(y_1,y_2)\big| \leq \sqrt2\,d_3(x,y).\]
This is a general principle rather than a computation: \textbf{Lipschitz continuity survives
passage to an equivalent metric}, with the constant multiplied by the comparison factor.
Continuity survives too, being topological, and by \cref{ex:2.5} the three product metrics
induce the same topology.

% Reclassified from ainote to remark: sharpness of the constant is mathematics.
\begin{remark}[The constant $2$ cannot be improved]
The constant $2$ for $d_1$ is sharp, so $d$ is genuinely \emph{not} $1$-Lipschitz with respect
to $d_1$. Take $X = \mathbb{R}$ with $x := (0,0)$ and $y := (-1,1)$. Then $d(x_1,x_2) = 0$ and
$d(y_1,y_2) = 2$, while $d_1(x,y) = \max\{1,1\} = 1$, so the difference equals exactly
$2 = 2\,d_1(x,y)$. This is worth checking, because the problem asks only \emph{whether} $d$ is
Lipschitz for $d_1$ and $d_3$, and it is easy to answer \qt{yes, with the same constant} without
noticing that the constant has in fact changed.
\end{remark}
\end{exercisesolution}

\begin{exercisesolution}[Solution to \cref{ex:2.7}]
% Generator: Claude Opus 5 (High)
All three arguments are short; the point of giving more than one is to see how differently the
same fact reads through the three definitions of \cref{def:continuous}.

\textbf{Topological (\cref{def:continuous}\textbf{(c)}).} This is the shortest, because
preimages compose. Let $W \subseteq Z$ be open. Then
\[(g\circ f)^{-1}(W) = f^{-1}\big(g^{-1}(W)\big),\]
and $g^{-1}(W)$ is open in $Y$ by continuity of $g$, so its preimage under $f$ is open in $X$ by
continuity of $f$. Hence $g\circ f$ is continuous.

\textbf{Sequential (\cref{def:continuous}\textbf{(b)}).} Let $x_n \to x$ in $X$. Continuity of
$f$ gives $f(x_n) \to f(x)$ in $Y$, and continuity of $g$ applied to \emph{that} sequence gives
$g(f(x_n)) \to g(f(x))$ in $Z$.

\textbf{$\varepsilon$--$\delta$ (\cref{def:continuous}\textbf{(a)}).} Fix $x_0 \in X$ and
$\varepsilon > 0$. Continuity of $g$ at $f(x_0)$ supplies $\eta > 0$ such that
$d_Y(f(x_0),y) < \eta$ implies $d_Z\big(g(f(x_0)),g(y)\big) < \varepsilon$. Continuity of $f$ at
$x_0$, applied to \emph{that} $\eta$, supplies $\delta > 0$ such that $d_X(x_0,x) < \delta$
implies $d_Y(f(x_0),f(x)) < \eta$. Chaining the two, $d_X(x_0,x) < \delta$ implies
$d_Z\big((g\circ f)(x_0),(g\circ f)(x)\big) < \varepsilon$.

\begin{remark}[Why the topological proof is the short one]
All three arguments have the same shape, applying the hypothesis for $g$ first and then feeding
the result to $f$. Only the topological version needs no quantifiers at all, because the set
identity $(g\circ f)^{-1} = f^{-1}\circ g^{-1}$ does the bookkeeping the other two proofs carry
out by hand. This is a fair advertisement for \cref{def:continuous}\textbf{(c)}, which can look
like the least intuitive of the three on first meeting: it is the formulation in which
composition is trivial, and that is exactly why general topology takes it as \emph{the}
definition.
\end{remark}
\end{exercisesolution}

\begin{exercisesolution}[Solution to \cref{ex:continuity_classification}]
% Generator: Claude Sonnet 5 (Medium)
\begin{itemize}
  \item $\sqrt{x}$ on $[0,\infty)$: \textbf{uniformly continuous, not Lipschitz.} Uniform
    continuity follows from continuity on $[0,1]$ (compact, \cref{item:uniform_continuity_compact})
    together with $\tfrac12$-Lipschitz behaviour on $[1,\infty)$ (since $|\sqrt x - \sqrt y| \leq
    \tfrac12|x-y|$ there by the mean value theorem, as $|(\sqrt{\cdot})'| = \tfrac{1}{2\sqrt x}
    \leq \tfrac12$ for $x \geq 1$). Not Lipschitz near $0$: for $x_n :=
    1/n^2 \to 0$, $\dfrac{|\sqrt{x_n}-\sqrt0|}{|x_n-0|} = n \to \infty$, so no global $L$ works.
  \item $1/x$ on $(0,1)$: \textbf{continuous, neither uniformly nor Lipschitz continuous.} For
    $x_n := 1/n$, $y_n := 1/(2n)$, $|x_n-y_n| \to 0$ but $|1/x_n - 1/y_n| = n \to \infty$, so no
    uniform $\delta(\varepsilon)$ (and hence no Lipschitz constant) can work near $0$.
  \item $\sgn(x)$ on $(-1,1)$: \textbf{none of the three}, being discontinuous at $0$.
  \item $|x|^{3/2}$ on $(-1,1)$: \textbf{Lipschitz, hence also uniformly continuous and
    continuous.} $|x|^{3/2}$ is $C^1$ on $(-1,1)\setminus\{0\}$ with $\big|\tfrac{d}{dx}|x|^{3/2}\big|
    = \tfrac32|x|^{1/2} \leq \tfrac32$, and the bound extends across $0$ by direct estimation, so
    $L := \tfrac32$ works globally.
\end{itemize}
\end{exercisesolution}

\begin{exercisesolution}[Solution to \cref{ex:3.1}]
% Generator: Claude Opus 5 (High)
Both parts ask for a map with no fixed point, and each defeats one hypothesis of
\cref{thm:banach_fixed_point} while keeping the other, which is exactly what
\cref{rem:banach_hypotheses} predicts must be possible.
\begin{enumerate}[label=\textbf{(\alph*)}]
  \item Take $f(x) := \tfrac{x+1}{2}$. It maps $[0,1)$ into $[\tfrac12,1) \subseteq [0,1)$, and
    $|f(x)-f(y)| = \tfrac12|x-y|$, so it is $\tfrac12$-Lipschitz. A fixed point would satisfy
    $x = \tfrac{x+1}{2}$, i.e.\ $x = 1$, which is not in $[0,1)$. \textbf{What fails is
    completeness:} $[0,1)$ is not a complete metric space, the iterates
    $x_{n+1} = f(x_n)$ march towards the missing endpoint $1$, and \cref{thm:banach_fixed_point}
    does not apply.
  \item Take the translation $f(x) := x + (1,0)$. It is an isometry, since
    $|f(x)-f(x')| = |x-x'|$ exactly, and $f(x) = x$ would force $(1,0) = 0$. \textbf{What
    fails is the strict contraction:} an isometry has Lipschitz constant $L = 1$, not $L < 1$,
    and $\mathbb{R}^2$ is perfectly complete. Together with part \textbf{(a)} this shows neither
    hypothesis of the fixed point theorem can be dropped.
\end{enumerate}
\end{exercisesolution}

\begin{exercisesolution}[Solution to \cref{ex:3.2}]
% Generator: Claude Opus 5 (High)
\textbf{No continuous extension exists for $xy/(x^2+y^2)$.} Suppose $f$ were continuous on
$\mathbb{R}^2$ and agreed with $xy/(x^2+y^2)$ off the origin. Approach the origin along two
different lines:
\[f\!\left(\tfrac1k,\tfrac1k\right) = \frac{1/k^2}{2/k^2} = \frac12,
  \qquad\qquad
  f\!\left(\tfrac1k,0\right) = 0 \quad\text{for every } k.\]
Both sequences tend to $(0,0)$, so continuity would force
$f(0,0) = \tfrac12$ and $f(0,0) = 0$ simultaneously. Hence no such $f$ exists. (This is the same
function, and the same argument, as \cref{ex:partial_derivatives_not_continuous},
where it reappears to show that the existence of both partial derivatives implies nothing.)

\textbf{Exactly one continuous extension exists for $xy/\sqrt{x^2+y^2}$.} Write
$(x,y) = (r\cos\theta, r\sin\theta)$ with $r > 0$. Then
\[g(x,y) = \frac{r^2\cos\theta\sin\theta}{r} = r\cos\theta\sin\theta,
  \qquad\text{so}\qquad |g(x,y)| \leq \frac r2 = \frac{|(x,y)|}{2}.\]
The bound is uniform in $\theta$, so $g(x,y) \to 0$ as $(x,y) \to (0,0)$, and setting
$g(0,0) := 0$ produces a continuous function on all of $\mathbb{R}^2$.

\textit{Uniqueness} is the part worth spelling out, and it needs no computation at all: the value
at the origin is \emph{forced}. Since $(0,0)$ lies in the closure of
$\mathbb{R}^2\setminus\{(0,0)\}$, any two continuous extensions agree on a dense subset of
$\mathbb{R}^2$, and two continuous functions agreeing on a dense set agree everywhere. To see the
latter, take a sequence from the dense set converging to the point in question and apply
\cref{def:continuous}\textbf{(b)} to both. So the extension is unique whenever it exists, which
is also why the first half of the problem had only to exhibit \emph{two} conflicting limits.
\end{exercisesolution}

\begin{exercisesolution}[Solution to \cref{ex:3.6}]
% Generator: Claude Opus 5 (High)
\textbf{(a), (c) and (d) are true; only (b) is false.}

\textbf{(a) True.} Addition is Lipschitz, hence uniformly continuous
(\cref{rem:continuity_hierarchy}):
\[\big|(x+y)-(x'+y')\big| \leq |x-x'| + |y-y'| \leq \sqrt2\,\big|(x,y)-(x',y')\big|,\]
the last step by Cauchy--Schwarz. So $\delta := \varepsilon/\sqrt2$ works everywhere at once.

\textbf{(b) False.} Multiplication is continuous but not uniformly so on the unbounded domain
$\mathbb{R}^2$. Take
\[p_k := \left(k,\ \tfrac1k\right), \qquad q_k := \left(k,\ 0\right).\]
Then $|p_k - q_k| = \tfrac1k \to 0$, while the values are $k \cdot \tfrac1k = 1$ and
$0$, a constant gap of $1$. So no $\delta$ can work for $\varepsilon < 1$. The failure is exactly
the one described in \cref{rem:continuity_vs_uniform_continuity}: the function steepens without
bound as one runs off to infinity, so the local $\delta$ shrinks without a positive floor.

\textbf{(c) and (d) True,} and both for free: $[0,1]^2$ is closed and bounded, hence
\textbf{compact} by \cref{thm:heine_borel}, and a continuous function on a compact metric space
is automatically uniformly continuous by
\cref{prop:compact_implies_uniformly_continuous} (Corsin's third reason to care about
compactness, \cref{item:uniform_continuity_compact}). No estimate is needed for either.

\begin{remark}
The pair \textbf{(b)} versus \textbf{(d)} is the whole lesson: the \emph{same} function $xy$ is
uniformly continuous on $[0,1]^2$ and not on $\mathbb{R}^2$. Uniform continuity is therefore not
a property of a formula but of a formula \emph{together with its domain}, and compactness of the
domain is what removes the room to run off to a bad point.
\end{remark}
\end{exercisesolution}

\begin{exercisesolution}[Solution to \cref{ex:3.11}]
% Generator: Claude Opus 5 (High)
\begin{enumerate}[label=\textbf{\arabic*.}]
  \item \textbf{($\Leftarrow$) A modulus of continuity forces uniform continuity.} Suppose
    $d_Y(f(x),f(y)) \leq \omega(d_X(x,y))$ for all $x,y$, with $\omega$ as in the statement. Let
    $\varepsilon > 0$. Since $\omega(t) \to 0$ as $t \downarrow 0$, there is $\delta > 0$ with
    $\omega(\delta) < \varepsilon$. If $d_X(x,y) < \delta$ then, $\omega$ being nondecreasing,
    \[d_Y\big(f(x),f(y)\big) \leq \omega\big(d_X(x,y)\big) \leq \omega(\delta) < \varepsilon.\]
    The $\delta$ did not depend on $x$ or $y$, which is precisely
    \cref{def:uniformly_continuous}.

    \textbf{($\Rightarrow$) Uniform continuity produces one.} Following the official hint, define
    the \newterm{modulus of continuity of $f$} as its worst oscillation at scale $t$:
    \[\omega_f(t) := \sup\big\{\,d_Y(f(x),f(y))\ :\ x,y \in X,\ d_X(x,y) \leq t\,\big\}
      \ \in\ [0,\infty].\]
    It has all three required properties. $\omega_f(0) = 0$, because $d_X(x,y) \leq 0$ forces
    $x=y$ by definiteness and then $d_Y(f(x),f(y)) = 0$. It is nondecreasing, because enlarging
    $t$ enlarges the set over which the supremum is taken. And $\omega_f(t) \to 0$ as
    $t \downarrow 0$: given $\varepsilon > 0$, uniform continuity supplies $\delta > 0$ with
    $d_X(x,y) < \delta \Rightarrow d_Y(f(x),f(y)) < \varepsilon$, so every $t < \delta$ has
    $\omega_f(t) \leq \varepsilon$. Finally $f$ has modulus $\omega_f$ by construction, since
    $d_Y(f(x),f(y))$ is one of the numbers the supremum defining $\omega_f(d_X(x,y))$ ranges
    over.
  \item \textbf{(Bonus 1) On an interval, $\omega_f$ is already subadditive.} Let
    $t_1,t_2 \geq 0$ and take any $x,y \in X$ with $|x-y| \leq t_1+t_2$. Choose an intermediate
    point $z$ on the segment from $x$ to $y$ as follows: if $|x-y| \leq t_1$ put $z := y$;
    otherwise let $z$ be the point of that segment at distance exactly $t_1$ from $x$, so that
    $|z-y| = |x-y| - t_1 \leq t_2$. Either way $z \in X$, because $X$ is an interval and $z$ lies
    between $x$ and $y$, and $|x-z| \leq t_1$, $|z-y| \leq t_2$. The triangle inequality in $Y$
    now gives
    \[d_Y\big(f(x),f(y)\big) \ \leq\ d_Y\big(f(x),f(z)\big) + d_Y\big(f(z),f(y)\big)
      \ \leq\ \omega_f(t_1) + \omega_f(t_2).\]
    Taking the supremum over all admissible $x,y$ yields
    $\omega_f(t_1+t_2) \leq \omega_f(t_1) + \omega_f(t_2)$.
  \item \textbf{(Bonus 2) Convexity is exactly what the argument needed.} Repeat part 2 verbatim,
    now reading $|x-y|$ as the Euclidean norm on $\mathbb{R}^n$ rather than the absolute value on
    $\mathbb{R}$: the only property of an interval that was used is that the segment from $x$ to
    $y$ lies inside $X$, and that is the definition of convexity. The intermediate point is
    $z := x + \min\{t_1,|x-y|\}\,\frac{y-x}{|y-x|}$ for $x \neq y$, which lies on the segment and
    satisfies the same two estimates.
\end{enumerate}

\begin{remark}[Why the value $+\infty$ has to be allowed]
The official remark accompanying this problem gives the reason, and it is worth restating: on
$X = \mathbb{Z} \subseteq \mathbb{R}$ the map $f(n) := n^2$ is uniformly continuous for the
trivial reason that any two distinct points are at distance at least $1$. Take
$\delta := \tfrac12$, and the implication holds vacuously. Yet
$|f(n+1)-f(n)| = 2n+1$ is unbounded, so $\omega_f(1) = +\infty$. Parts 2 and 3 rule this out by
hypothesis rather than by luck: on a convex set the intermediate-point construction lets a large
displacement be broken into many small ones, so $\omega_f(t) \leq \lceil t/s\rceil\,\omega_f(s)$
is finite as soon as some $\omega_f(s)$ is. A space like $\mathbb{Z}$ has no intermediate points
to subdivide with.
\end{remark}
\end{exercisesolution}
